\section{Introduction}\label{sec:introduction}

Land cover data products provide valuable information about the characteristics of the Earth's surface for a wide variety of applications, including agriculture, urban planning, and environmental monitoring~\citep{scarpace1980land, lovelandDevelopmentGlobalLand2000,townshendLandCover1992,wulderLandCover2018,banGlobalLandCover2015}.
Research into automated methods for land cover classification from remotely sensed data has been ongoing for several decades, with a majority of methods relying on supervised learning techniques to classify land cover using multispectral aerial or satellite imagery~\citep{brownLessonsLearnedImplementing2020, brownDynamicWorldRealtime2022, friedlGlobalLandCover2002, friedlMODISCollection52010, zhuContinuousChangeDetection2014}. 
Historically, pixel-wise and object-based image analysis (OBIA) classification methods have been the predominant approach for land cover classification, where spatial units (i.e., pixels or objects) in the image are classified based on their spectral reflectance, hand-crafted features, and ancillary data from other geospatial datasets~\citep{mciverEstimatingPixelscaleLand2001, khatamiMetaanalysisRemoteSensing2016}. 
Very high spatial resolution (VHSR, \(<\) \SI{5}{\meter}) land cover products offer more detailed information about the Earth's surface with greater spatial accuracy, which is valuable for applications that require fine-grained information about land cover types and their spatial distribution \citep{liImpactsSpatialResolutions2022,fisherImpactSatelliteImagery2018}.
However, the creation of such VHSR data products is inherently challenging due to the complexity of the land cover patterns at these scales and the increased spectral variability within classes~\citep{strahlerNatureModelsRemote1986,woodcockFactorScaleRemote1987,hansenReviewLargeArea2012}.
In this context, the application of deep convolutional neural networks (CNNs) to land cover classification tasks has gained significant attention in recent years, as these models excel at extracting spatial and spectral features from imagery in a data-driven manner, making them well-suited for VHSR land cover classification tasks~\citep{zhuDeepLearningRemote2017a, maDeepLearningRemote2019, valiDeepLearningLand2020}.
Nevertheless, these models require far more labeled training data compared to traditional machine learning methods, serving as a significant barrier to their widespread adoption for operational land cover mapping at VHSR~\citep{zhaoReviewConvolutionalNeural2024, huTransferringDeepConvolutional2015}.

% Object-based Image Analysis (OBIA) techniques have been proposed as a potential workflow to address these issues, but their performance is highly dependent on the quality of objects or superpixels extracted from the imagery during the segmentation process, which must be carefully tuned to the specific characteristics of the imagery and the land cover classes of interest~\citep{blaschkeGeographicObjectBasedImage2014a, johnsonUnsupervisedImageSegmentation2011, maReviewSupervisedObjectbased2017}.
% This pixel-wise paradigm is effective for classifying coarse and moderate spatial resolution imagery (e.g., MODIS, Landsat, Sentinel-2), but is less effective for classifying high spatial resolution imagery due to decreased inter-class spectral separability and increased intra-class spectral variability~\citep{strahlerNatureModelsRemote1986,woodcockFactorScaleRemote1987,hansenReviewLargeArea2012, grekousisOverview21Global2015, markhamLandCoverClassification1981, luSurveyImageClassification2007}. Object-based Image Analysis (OBIA) techniques have been proposed as a potential workflow to address these issues, but their performance is highly dependent on the quality of objects or superpixels extracted from the imagery during the segmentation process, which must be carefully tuned to the specific characteristics of the imagery and the land cover classes of interest~\citep{blaschkeGeographicObjectBasedImage2014a, johnsonUnsupervisedImageSegmentation2011, maReviewSupervisedObjectbased2017}.

% The application of deep computer vision models to land cover classification has been of significant interest in recent years, as these models excel at extracting spatial and spectral features from imagery in a data-driven manner, making them well-suited for high spatial resolution land cover classification tasks~\citep{zhuDeepLearningRemote2017a, maDeepLearningRemote2019, valiDeepLearningLand2020}.

% CNNs are a class of deep learning models that can automatically extract hierarchical spatial structures from the input imagery that are useful for classifying land cover type without extensive feature engineering~\citep{liSurveyConvolutionalNeural2022a, songPatchBasedLightConvolutional2019, panLandcoverClassificationMultispectral2020,kussulDeepLearningClassification2017}.
% CNN-based semantic segmentation architectures learn a mapping from the input image to a pixel-wise classification probability map, which can be used to generate a land cover map by assigning each pixel to the class with the highest probability~\citep{longFullyConvolutionalNetworks2015}.
% CNN-driven semantic segmentation models have been widely used for VHSR land cover classification tasks, as they leverage both spectral and spatial information to classify land cover in the input imagery under a unified, single-step framework~\citep{yuanReviewDeepLearning2021}.
% However, these models require large amounts of labeled training data to learn the complex spatial and spectral patterns present in the imagery, which can be difficult to obtain for high spatial resolution imagery due to the high cost of collecting and labeling the data~\citep{huTransferringDeepConvolutional2015}.

This need for large amounts of ground-truth data for accurate training of deep learning models has led to extensive research into methods for reducing the amount of data necessary to train these models~\citep{wurmSemanticSegmentationSlums2019,tuiaDomainAdaptationClassification2016, martinsDeepLearningHigh2022}.
Most of these methods leverage transfer learning procedures, where the parameters of a model pre-trained using a large dataset are used to initialize a model that is subsequently fine-tuned on a smaller dataset that is more relevant to the target task~\citep{weissSurveyTransferLearning2016, zhuangComprehensiveSurveyTransfer2021}.
For example, a common approach to transfer learning for land cover classification tasks is to pre-train a model on a large general-purpose image classification dataset (e.g., ImageNet~\citep{russakovskyImageNetLargeScale2015}) and then fine-tune the model on a smaller dataset of interest~\citep{piresdelimaConvolutionalNeuralNetwork2020,marmanisDeepLearningEarth2016,nogueiraBetterExploitingConvolutional2017, zouTransferLearningClassification2018}. 
Despite being effective, the domain shift between the source and target datasets can hinder the model's performance, and the extent of this domain shift between general-purpose image classification datasets and high spatial resolution remote sensing datasets is not well understood~\citep{maTransferLearningEnvironmental2024}.
As fully-supervised pre-training requires large-scale labeled datasets, which are relatively scarce in the remote sensing domain for VHSR imagery, there is a need for alternative pre-training strategies that reduce the reliance on labeled data~\citep{wangEmpiricalStudyRemote2023, xuSelfsupervisedPretrainingLargescale2024}.
% \citep{bastaniSatlasPretrainLargeScaleDataset2023,manasSeasonalContrastUnsupervised2021a,jeanTile2VecUnsupervisedRepresentation2019}.

% These approaches are rapidly gaining popularity as a means to pre-train large models using remote sensing imagery that can serve as strong feature extraction backbones for a variety of downstream tasks in remote sensing analysis without the need for ground truth data.


% As part of a larger movement in the machine learning community, self-supervised learning techniques have been proposed as a more flexible approach for building large deep learning models that produce high-quality feature representations without the need for labeled data \citep{liuSelfSupervisedLearningGenerative2023}.
More recently, self-supervised pre-training methods have been proposed as an alternative to traditional fully-supervised pre-training methods, enabling models to learn useful spatial and spectral information from input imagery through the use of pretext tasks that do not require labeled data~\citep{wuUnsupervisedFeatureLearning2018,tianContrastiveMultiviewCoding2020, chenSimpleFrameworkContrastive2020}.
These techniques are particularly attractive for remote sensing applications, as large amounts of labeled imagery datasets are often unavailable, but source imagery is widely available due to the growing number of Earth observation satellites and aerial imagery platforms~\citep{luVisionFoundationModels2025}.
However, few studies have rigorously investigated the applicability of these methods to operational land cover mapping at VHSR, especially in scenarios where ground-truth data are scarce.
% , though similar studies have shwon  promising results in medium resolution land cover classification tasks (\SIrange{10}{30}{\meter}) ~\citep{ayushGeographyAwareSelfSupervisedLearning2021, manasSeasonalContrastUnsupervised2021a}.
Most research evaluating self-supervised models pre-trained with remote sensing data uses benchmark VHSR land cover datasets as a proxy to evaluate the quality of pre-trained models, with the ISPRS Potsdam~\citep{2DSemanticLabeling} and Vaihingen~\citep{2DSemanticLabel} datasets seeing wide use for this purpose~\citep{sunRingMoRemoteSensing2023, wangMTPAdvancingRemote2024, guoSkySenseMultiModalRemote2024, wangAdvancingPlainVision2023}.
However, these datasets contain relatively large amounts of labeled training data, and thus may not necessarily be representative of real-world scenarios where no available VHSR land cover data exists, and costly manual annotation (e.g., by experience, a single 256×256 \SI{1}{\meter} image patch can take upwards of 30 minutes to label well) is required to create this data.
% Given the growing interest in self-supervised techniques for building large foundation models for remote sensing applications~
Furthermore, while there is growing momentum in the development of large remote sensing foundation models using self-supervised techniques \citep{xiaoFoundationModelsRemote2025,luVisionFoundationModels2025}, there is still a critical need to systematically assess how these representations can enhance VHSR land cover mapping performance in practical contexts, a gap this study directly addresses.
% further investigation into the practical utility of these self-supervised learning techniques for operational land cover mapping at VHSR under extreme label-scarce scenarios is needed.
% , such as the ISPRS Potsdam and Vaihingen datasets~\citep{isprs2012isprs}, which are relatively small and may not be representative of the challenges associated with operational land cover mapping at VHSR over large areas.

% One such approach to self-supervised learning uses contrastive learning: these contrastive methods train a model to produce similar feature representations for different augmented views of the same input while enforcing dissimilarity between feature representations of different inputs \citep{jaiswalSurveyContrastiveSelfSupervised2021}.
% However, contrastive methods require large batch sizes due to the need for many negative samples to enforce dissimilarity between different inputs, which can induce substantial computational overhead during training and limit the size of models that can be trained \citep{chenSimpleFrameworkContrastive2020, chenWhyWeNeed2022}.
Self-supervised approaches inspired by masked language modeling in natural language processing have been shown to be strong approaches for learning useful feature representations from images without the need for labeled data~\citep{heMaskedAutoencodersAre2022, xieSimMIMSimpleFramework2022}.
Masked image modeling (MIM) methods take advantage of the Vision Transformer (ViT) patching mechanism to randomly mask out a portion of the input image, then train the model to reconstruct the missing patches from the visible patches~\citep{linSSMAESpatialSpectral2023,wangFeatureGuidedMasked2025}.
MIM techniques have become the predominant self-supervised approaches in remote sensing for creating large foundation models, which are large pre-trained models that can be adapted for a wide variety of tasks and domains that ideally outperform models trained from scratch or use more traditional transfer learning approaches~\citep{hongSpectralGPTSpectralRemote2024,wangAdvancingPlainVision2023}.
% Still, most MIM techniques require ViT backbones, as they rely on the patching mechanism and global feature extraction capacity of attention mechanisms.
However, while powerful, ViTs lack the inductive biases toward imagery that make CNNs effective at extracting spatial features from images~\citep{adadiSurveyDataefficientAlgorithms2021}.
This is problematic, as ViTs typically only begin to outperform CNNs when pre-trained on very large datasets~\citep{dosovitskiyImageWorth16x162020}, a practical consideration for researchers and practitioners who must balance both computational resources and data availability when developing models.
Self-distillation techniques, such as Bootstrap Your Own Latent (BYOL)~\citep{grillBootstrapYourOwn2020} and DINO~\citep{caronEmergingPropertiesSelfSupervised2021}, are alternatives to MIM pre-training that can be used to learn robest and transferable feature representations without the need for negative samples, large batch sizes, or ViT backbones, which are hard requirements for many other self-supervised learning methods~\citep{chenSimpleFrameworkContrastive2020, chenWhyWeNeed2022}.
Such approaches largely revolve around maximizing the similarity between feature representations of different augmented views of the same input using a teacher-student framework~\citep{heFoundationModelBasedMultimodal2024}.
% However, instead of relying on negative samples to prevent the network from learning trivial solutions, these methods utilize a teacher-student framework, where a student (or online) network is trained to predict the output of a teacher (or target) network that receives a different augmented view of the same input.
While the student network is trained directly using standard backpropagation and gradient descent algorithms, the teacher's parameters are updated using a momentum-based rule derived from the student's parameters.
Self-distillation methods are attractive for remote sensing applications due to their inherent flexibility and ability to learn strong feature representations using a variety of backbone architectures, including CNNs~\citep{chenExploringSimpleSiamese2020}.
Given the architectural constraints of MIM and the computational demands of contrastive methods, self-distillation techniques offer a compelling, architecture-agnostic solution, suitable for pre-training CNNs on domain-specific, moderate-sized datasets~\citep{jainSelfSupervisedLearningInvariant2022}.

In this study, we propose an innovative label-efficient framework that addresses the challenge of operational land cover mapping at \SI{1}{\meter} resolution under label-scarce scenarios through self-supervised learning.
Specifically, our framework uses BYOL as a pretext task during pre-training, enabling a ResNet-101 convolutional encoder network to learn to extract deep features from unlabeled 377,921 256\(\times\)256 patches of \SI{1}{\meter}  color-infrared aerial imagery without the need for full supervision using labeled data.
We then transfer the learned encoder weights into multiple semantic segmentation architectures (FCN, U-Net, Attention U-Net, DeepLabV3+, UPerNet, PAN) that are subsequently fine-tuned using very small training dataset sizes (250, 500, and 750 256\(\times\)256 patches) using cross-validation.
We classify 8 classes of land cover (open water, impervious structures, impervious surfaces, barren land, forest/woody vegetation, herbaceous/low vegetation, cultivated crops, and unclassified (i.e., shadow)) over the entire state of Mississippi, USA using an ensemble of the best-performing models, totaling over 123 billion \SI{1}{\meter} pixels.
% This approach is used to tackle land cover classification at \SI{1}{\meter} resolution over the state of Mississippi using only 1,000 256\(\times\)256 patches of ground truth data. 
Validation and accuracy assessment are performed using a spatially independent test set of 25,000 randomly sampled points with manually verified labels to ensure robust and unbiased evaluation of our approach and the resulting land cover map.
We demonstrated that these approaches can be used to train deep semantic segmentation models that are accurate, generalizable, and capable of yielding reasonable performance even under extreme low-data scenarios, even when only 3 bands (red, green, and near-infrared) are available in the input imagery.
% The contributions of this work are as follows:
% (1) we showed that self-supervised learning is an effective strategy for pre-training convolutional encoders using VHSR aerial imagery without the need for labeled data;
% (2) we demonstrated that these learned feature representations can be successfully transferred to multiple semantic segmentation architectures and fine-tuned using very small training dataset sizes to achieve satisfactory performance for \SI{1}{\meter} land cover classification; and
% (3) we produce a comprehensive \SI{1}{\meter}, 8-class land cover map over the state of Mississippi, USA using only 1,000 256\(\times\)256 patches of ground truth data, the first such map of its kind produced using such a small amount of training data over a large area.
Our ultimate aim is to reduce practical barriers to producing high spatial resolution land cover maps while contributing to ongoing efforts to evaluate and contextualize the implications of self-supervised learning techniques for remote sensing purposes.
% This research contributes to ongoing efforts to evaluate and contextualize the implications of self-supervised learning techniques for remote sensing purposes by proving a practical application of these methods to a real-world land cover classification task at VHSR.
